% Generated draft: requires mathematical and source review.
\documentclass[11pt]{article}
\usepackage[margin=25mm]{geometry}
\usepackage{fontspec}
\setmainfont{Latin Modern Roman}
\usepackage{amsmath,amssymb,mathtools,unicode-math}
\setmathfont{Latin Modern Math}
\usepackage{enumitem,needspace,xurl,hyperref}
\setlength{\parindent}{0pt}
\setlength{\parskip}{5pt}
\setlength{\emergencystretch}{3em}
\newcommand{\N}{\mathbb{N}}
\newcommand{\Z}{\mathbb{Z}}
\newcommand{\Q}{\mathbb{Q}}
\newcommand{\R}{\mathbb{R}}
\newcommand{\C}{\mathbb{C}}
\newcommand{\F}{\mathbb{F}}
\newcommand{\A}{\mathbb{A}}
\newcommand{\PP}{\mathbb{P}}
\newcommand{\E}{\mathbb{E}}
\newcommand{\eps}{\varepsilon}
\newcommand{\dd}{\,\mathrm{d}}
\begin{document}
\section*{UnsolvedMath ID: 30004668}
\textbf{Source problem code:} catalogue entry
\subsection{Definitions and mathematical statement}
Let \(U^+\), \(K(Q,W)\), and \(\psi:U^+\to K(Q,W)\) denote the specific graded algebras and graded algebra homomorphism intended in the cited Oberwolfach report. Define
\[
\ker(\psi)=\{x\in U^+:\psi(x)=0_{K(Q,W)}\}.
\]
The supplied conjecture is
\[
\ker(\psi)=\{0_{U^+}\},
\]
or equivalently,
\[
\forall x,y\in U^+,\qquad \psi(x)=\psi(y)\Longrightarrow x=y.
\]
This concerns the particular homomorphism in the report. The supplied record does not define \(U^+\), \(K(Q,W)\), \(Q\), \(W\), their coefficient rings or grading, or the construction of \(\psi\). Consequently, the injectivity target is explicit, but its mathematical objects cannot be specified self-containedly from the supplied evidence. Any hypotheses restricting these objects remain unresolved.
\subsection{Short English statement}
The conjecture asks whether the particular graded algebra homomorphism \(\psi:U^+\to K(Q,W)\) identifies no two distinct elements. Definitions of the algebras, the map, and their hypotheses are missing from the supplied record and require human mathematical review.
\subsection{Sources}
\begin{itemize}
\item Local registry: \texttt{data/unsolvedmath/problems.json}, permanent UnsolvedMath ID \(30004668\), source problem code \texttt{catalogue entry}, title ``Injectivity of the Quantum K-Theoretic Hall Homomorphism.'' The supplied collection is Oberwolfach Reports Open Problems; the proposed year is 2021, and no individual proposer is named.
\item Supplied report citation: \emph{Algebraic Groups} (2022), Oberwolfach Reports, DOI 10.4171/owr/2021/20, pp.\ 1087--1148. \url{https://doi.org/10.4171/owr/2021/20}. The supplied background describes Oberwolfach Reports as documenting mathematical workshops and their open-problem discussions. The citation's 2022 label and its 2021 report identifiers are retained without resolving their relationship.
\item Supplied literature source: Oberwolfach Report 20/2021. \url{https://oa.tib.eu/renate/server/api/core/bitstreams/f2f10f3f-72af-4586-9653-7e0b6b7e77c9/content}.
\item Supplied literature source: P\u{a}durariu, \emph{K-theoretic Hall algebras}. \url{https://arxiv.org/abs/1911.05526}.
\item Literature assessment dated exactly \texttt{2026-08-22}: the supplied curation labels the problem open, states that injectivity is explicitly conjectured in the report, and says that subsequent K-theoretic Hall-algebra literature still describes recovery of the quantum affine algebra as conjectural rather than proving this map injective. This is supplied dated curation derived from the OWR curation, not an independently verified assessment; its time-sensitive status and cited claims require independent verification.
\end{itemize}
\end{document}
